\section{Discussion}\label{sec:discussion}

\subsection{Where CKG Wins and Why}\label{sec:disc-ckg-wins}

% TODO: Expand with experimental evidence

CKG's advantages are structural:
\begin{itemize}
  \item \textbf{T2/T3 queries}: Explicit edges eliminate multi-hop inference errors.
    RAG must infer transitive dependencies from unstructured text; CKG traverses
    them directly via BFS/DFS.
  \item \textbf{T4 queries}: Taxonomy filtering achieves BC $\approx$ 1.0 by
    construction, as the TaxonomyID field provides exact category membership.
  \item \textbf{Hallucination}: HR $=$ 0 because CKG only returns concepts
    present in the source CSV. No generative step can introduce phantom entities.
  \item \textbf{RDS}: Near-zero build cost combined with 150--400 tokens per
    query yields order-of-magnitude efficiency gains.
\end{itemize}

\subsection{Where RAG Is Competitive}\label{sec:disc-rag}

RAG remains competitive in specific scenarios:
\begin{itemize}
  \item T1 entity lookup on large open-domain corpora where rich context aids
    natural language generation.
  \item Domains without stable taxonomy (rapidly evolving fields).
  \item When CKG construction cost exceeds the efficiency savings.
\end{itemize}

\subsection{GraphRAG's Position}\label{sec:disc-graphrag}

GraphRAG occupies a middle ground: better than RAG on multi-hop reasoning
(graph structure helps) but worse than CKG (dynamic extraction introduces
noise and hallucinated edges). GraphRAG is the most expensive system
(high build cost + high query cost). Its best use case is unstructured
corpora with no available expert taxonomy.

\subsection{The Structure Premium}\label{sec:disc-premium}

% TODO: Report correlation results

We hypothesize that the token efficiency gap between CKG and RAG is
proportional to the structural richness of the domain's DAG:
\begin{equation}
  \text{dag\_richness}(d) = \frac{\text{edges}}{\text{concepts}} \times
    \text{mean\_indegree} \times \frac{1}{\text{orphan\_rate}}
\end{equation}

If the correlation between structure\_premium and dag\_richness exceeds
$r = 0.7$, this validates the core theoretical claim: explicit structure is
the source of the efficiency gain, not domain selection bias.

\subsection{Limitations}\label{sec:disc-limitations}

\begin{itemize}
  \item \textbf{Structural query scope:} Ground truth for T2, T3, and T4 queries
    is derived directly from DAG edges. CKG retrieves this ground truth by
    construction; RAG and GraphRAG must infer it from text. The comparison
    therefore demonstrates that \emph{explicit structure outperforms inferred
    structure on structural tasks}---not that CKG is a general-purpose
    retrieval system. CKG is not designed to compete with RAG on explanatory
    queries.
  \item \textbf{T1 as boundary test:} CKG scores F1 $\approx$ 0.09 on T1
    entity lookup queries because the DAG contains no explanatory prose. For
    open-ended definitional knowledge retrieval, RAG remains the appropriate
    architecture. The low T1 score confirms the benchmark is not rigged: a
    system that only reads DAG edges appropriately fails at questions the DAG
    cannot answer.
  \item The McCreary corpus is educational---results may not generalize to
    legal, financial, or medical domains.
  \item CKG requires upfront expert curation investment; this cost is not
    reflected in the benchmark's zero-build-cost assumption.
  \item Ground truth derived from DAG edges may not capture all valid
    natural language answers.
  \item All systems use the same LLM (Claude Sonnet 4.6); results may
    differ with other models.
\end{itemize}
